\section{问题三：固定发布时刻的条件预报与真实触发识别}
\subsection{任务目标与可识别范围}
题目要求给出触发时间的95\%置信区间、预警时间、区间覆盖率和空间距离。现有资料只有一个口岸冲击时刻，没有源区冰体失稳、物料释放或泥石流形成的独立时间观测；附件4的表面速度是2014--2024气候态栅格，没有2026年灾前时间维。因而现有资料不足以识别真实释放时刻$t_r$、形成时刻$t_f$、实际到达$t_a$及其预警和覆盖率。

本文把时间对象分开。$t_m$为源区融水条件首次满足的代理时刻，$t_a^{(0)}$为“融水出现即释放既定物料”情景下的口岸条件到达，$t_a$为实际到达。三者满足
\begin{equation}
t_r\ne t_m,\qquad
t_f=t_r+\Delta_{r\to f},\qquad
t_a=t_r+\Delta_{r\to\mathrm{inj}}+\Delta_{\mathrm{inj}\to a},
\label{eq:times}
\end{equation}
其中$\Delta_{r\to f}$为释放到流态转变的延迟，$\Delta_{r\to\mathrm{inj}}$为释放到公共注入断面的传播时间，$\Delta_{\mathrm{inj}\to a}$为注入到口岸的运行时间。没有释放观测时，本文不把$t_m$当作$t_r$，也不把$t_a^{(0)}$写成实际到达。表\ref{tab:q3-targets}列出各目标的处理。

\begin{table}[htbp]\centering\footnotesize
\caption{第三问目标与可识别性}\label{tab:q3-targets}
\begin{tabular}{@{}l l l@{}}\toprule
目标&可识别性&本文输出\\\midrule
$t_m$：源区融水起始代理&辅助气象目标&条件95\%区间\\
$t_r$：冰岩物料实际释放&不可识别&空值\\
$t_f$：泥石流形成&不可识别&空值\\
$t_a$：口岸实际到达&不可识别&空值\\
$t_a^{(0)}$：即时释放条件到达&辅助情景&条件95\%区间\\
10:30口岸冲击&题面单一锚点&事后单点包含\\
\bottomrule
\end{tabular}\end{table}

因此，本文给出的产品是$t_m$与$t_a^{(0)}$的条件区间、状态计数和空间距离。真实$t_r$、$t_f$、$t_a$、真实预警和覆盖率尚不能给出；当前结果属于辅助条件产品，不作为真实触发预报。

主发布时刻固定为2026年8月26日08:00当地时（00 UTC），预报窗口为
\begin{equation}
\mathcal H=[2026\text{-}08\text{-}26\,08{:}00,\,
2026\text{-}08\text{-}27\,08{:}00).
\end{equation}
主起报为8月25日18 UTC，另取同日06、12 UTC两次单次起报构造模式演化差异。三者按六小时分发上界分别于25日12 UTC、18 UTC和26日00 UTC可用，均不晚于发布时刻；每次起报生成1000个成员，合计$N=3000$。题面10:30只用于事后单点包含，不进入成员生成、系数拟合或配置选择。

\subsection{变量时间对齐与历史能量关系}
普通小时短波是前一小时均值，降水是前一小时累计，温湿度是对应时点的瞬时值。若按原始标签把三者放在同一行，会把前一小时间接辐射归入后一小时支持。本文按变量真实支持对齐到统一半开区间$[t,t+1)$：
\begin{equation}
SW_t=SW^{\mathrm{raw}}_{t+1},\qquad
P_t=P^{\mathrm{raw}}_{t+1},\qquad
\bar T_t=\frac{T^{\mathrm{raw}}_{t}+T^{\mathrm{raw}}_{t+1}}{2},\qquad
\overline{RH}_t=\frac{RH^{\mathrm{raw}}_{t}+RH^{\mathrm{raw}}_{t+1}}{2}.
\label{eq:support}
\end{equation}
前两者是过去小时均值与累计量，后两者是时点值的梯形均值近似。先取小时均值，再施加正部与门控，不交换非线性操作；任一端点缺失则该小时不可评估，不用下一条记录跨过时间缺口。发布前最后完成支持$[07{:}00,08{:}00)$取原始标签08:00的短波与降水，以及两个温湿端点；完整轨迹的末支持起点为原始末时点减一小时。若按起始标签对齐，会把前一小时间接辐射归入后一小时支持，故本文不采用该口径。

历史站点四类小时标签缺少起止或瞬时元数据。主假设为MET时点端点梯形均值、RADM/SOIL/FLUX起点标签小时均值；对MET三种处理与其余三族起点/终点组合共24种情形重建净辐射与消融能量目标、重拟合回归和AR(1)，各情形使用同一流程，未按事件结果挑选。24种情形给出相同的成员首达支持集合，说明辅助区间对时间标签假设不敏感；真实释放仍需独立观测。

在2021--2025年6--9月共同有效小时上拟合
\begin{equation}
\hat e_m=\beta_0+\beta_T T^++\beta_R R^++\beta_H RH+\varepsilon,
\qquad
T^+=\max(T,0),\quad R^+=\max(SW,0),
\end{equation}
其中$\hat e_m$与$R^+$的单位为W/m$^2$，$T$为$^\circ$C，$RH$为百分数。训练集共11986行，估计得
\begin{equation}
\beta=(-119.0252,\,-5.8526,\,0.4706,\,1.7154),
\end{equation}
残差一阶自回归参数$\phi=0.400328$、创新标准差$\sigma_\varepsilon=41.015695$ W/m$^2$，边际标准差44.715577 W/m$^2$，连续小时对11772个。预报时仅当$T>0$、$SW>20$ W/m$^2$且$\hat e_m>0$时判定潜在融水；任一必需输入缺失记不可评估，不以零值填补。小时水当量按$m_s=\hat e_m/334000\times3600$换算。

发布前最后完成支持$[07{:}00,08{:}00)$没有成员满足活动判据。发布后24小时窗内，3000个源区成员和3000个传播成员均判定为有效，有效比例$p_{\mathrm{fin}}=1$；无事件、发布时已活跃、右删失、不可评估、未到达及口岸左删失的成员计数均为0，状态合计一致。可行动质量
\begin{equation}
p_{\mathrm{actionable}}=p_{\mathrm{act}}+p_{\mathrm{fin}}=0+1=1
\end{equation}
只表示给定模型集合中需要发出源区融水提示的成员比例，不是物理发生概率。

\subsection{条件区间与支持包络}
每个成员的小时支持原样保留为$[L_i,U_i)$，宽度为1 h。为给出事件时刻分位，在固定随机流中为每个有限成员在支持内均匀抽取代表时刻
\begin{equation}
\tilde t_i=L_i+\tau_i(U_i-L_i),\qquad \tau_i\sim U[0,1),
\end{equation}
再取2.5\%与97.5\%分位。$t_m$与$t_a^{(0)}$由同一成员、同一传播配置配对生成：
\begin{equation}
\tilde t_{a,i}^{(0)}=\tilde t_{m,i}+\Delta t_{\mathrm{fall},i}
+\Delta t_{\mathrm{run}}^{(c_i)},
\end{equation}
其中$\Delta t_{\mathrm{fall}}$为源区质心至公共注入断面的无摩擦首波近似，两种高程语义下为29.50--29.67 s；$\Delta t_{\mathrm{run}}^{(c_i)}$取自16个去重D8配置之一。配置等权抽样只表示模型集合设计，不是现实配置概率。

条件区间见表\ref{tab:q3-event-intervals}，均为2026年8月26日UTC+8时刻，保留到毫秒。它们是有限成员分布内的名义95\%分位，不是经真实多事件频率校准的信赖区间。

\begin{table}[htbp]\centering\footnotesize
\caption{条件事件时刻区间}\label{tab:q3-event-intervals}
\begin{tabular}{@{}l c c c c@{}}\toprule
对象&2.5\%分位&中位&97.5\%分位&宽度/h\\\midrule
源区$t_m$&09:03:22.847&10:13:44.990&10:57:51.047&1.907833\\
即时释放$t_a^{(0)}$&09:35:26.068&10:46:31.992&11:31:59.410&1.942595\\
\bottomrule
\end{tabular}\end{table}

完整支持包络与支持端点分位不是同一对象。前者为有限成员全部半开支持的凸包
\begin{equation}
E_{\mathrm{support}}=[\min_i L_i,\max_i U_i),
\end{equation}
可能含空档，不是成员支持的精确保集；后者为
\begin{equation}
E_{0.95,\mathrm{endpoint}}=[q_{0.025}(L_i),q_{0.975}(U_i)],
\end{equation}
只是端点边际分位的收缩诊断。两者均不替代表\ref{tab:q3-event-intervals}的代表时刻事件区间，具体数值见表\ref{tab:q3-envelopes}。

\begin{table}[htbp]\centering\footnotesize
\caption{支持包络}\label{tab:q3-envelopes}
\begin{tabular}{@{}l c c c c@{}}\toprule
对象&包络类型&起点&终点&宽度/h\\\midrule
$t_m$&完整支持&$[$09:00:00.000&11:00:00.000$)$&2.000000\\
$t_m$&端点2.5\%/97.5\%&$[$09:00:00.000&11:00:00.000$]$&2.000000\\
$t_a^{(0)}$&完整支持&$[$09:18:03.728&11:45:13.334$)$&2.452668\\
$t_a^{(0)}$&端点2.5\%/97.5\%&$[$09:26:42.679&11:45:13.170$]$&2.308470\\
\bottomrule
\end{tabular}\end{table}

三类时间对象的分工见图\ref{fig:q3-time-products}；图中只比较代表事件分位、完整支持包络和支持端点包络，不把不同定义混成一个“95\%触发区间”。

\paperfig{figures/03_时间区间.png}{0.78}{三类时间对象}{fig:q3-time-products}

\subsection{单事件包含、提示时差与真实预警}
真实泥石流形成没有独立真值，真实触发覆盖率不可估计，本文不给出数值。题面只给出一个口岸冲击点$y_h=10{:}30$。在即时释放辅助情景下，条件区间09:35:26.068至11:31:59.410包含该点，报告为1/1。该1/1表明辅助情景与题面锚点不矛盾，不代表实际到达覆盖率。实际到达覆盖记为不可估计（释放模型未识别），区间、确定覆盖和可能覆盖尚未给出；长期频率需独立资料。

融水提示规则在发布时刻确定：若$p_{\mathrm{actionable}}\ge0.5$，则在08:00发布源区融水提示。本次$p_{\mathrm{actionable}}=1$，故提示时刻为08:00。相对10:30锚点的信息时差为2.5 h，这是“融水提示”相对既有冲击锚点的时间差，不是真实灾害预警提前量。真实灾害预警的首次报警、可执行报警和提前量尚不能给出。

辅助到达区间相对10:30的提前量为
\begin{equation}
\Delta_{a,\mathrm{range}}
=[10{:}30-U_{t_a^{(0)}},\,10{:}30-L_{t_a^{(0)}}]
=[-1.033169,\,+0.909426]\ \mathrm{h},
\end{equation}
区间跨过锚点，故2.5 h信息时差不构成命中提前量；源区代理中位距发布2.229 h，也不是真实释放提前量。图\ref{fig:q3-coverage}给出辅助单点包含和提示延迟对照；真实释放、形成、实际到达和真实灾害预警均标注为不可估计。

\paperfig{figures/03_区间覆盖率.png}{0.78}{辅助单点包含与提示时差}{fig:q3-coverage}

\subsection{空间距离与右删失运程}
按同一源区、公共注入断面和口岸代理断面分别计算距离。源区至注入断面的水平距离为2041.449 m；注入断面至口岸代理的中心线沿沟距离为20882.237 m；组合沿沟路线为22923.686 m。源区至口岸代理的平面直线距离为14771.715 m。沿沟里程与平面直线距离不可相加成一个“距离范围”。

16个去重D8配置的口岸首达时间范围为1054.226--2683.668 s，即17.57--44.73 min。全部配置的终态分析域里程均为26045.494670 m，且右删失标志为真；该值是演算分析域末端的右删失前缘，不是自然最大运程，也不是实测危险区边界。空间对象与状态见图\ref{fig:q3-distance}。

\paperfig{figures/03_空间距离.png}{0.80}{沿沟距离与D8到达}{fig:q3-distance}

\subsection{真实释放的识别条件}
式\eqref{eq:times}中的$t_r$不能由融水代理替代。可检验的观测约束路径是渐进破坏的逆速度模型。若取得发布前各时点$u_j$的独立位移$d_j$、相关误差与可用时刻，则对加速段构造
\begin{equation}
v_j=\frac{d_j-d_{j-1}}{u_j-u_{j-1}},\qquad
\dot v=A v^\alpha,
\end{equation}
当$\alpha=2$时
\begin{equation}
v^{-1}(u)=a+bu,\qquad t_r=-a/b\quad(b<0).
\end{equation}
它把可观测加速连接到渐进破坏的预测外推，但仍需检验稳定运动基线、窗口敏感性、测量误差、速度求差相关性和外推失效。附件4的速度场没有2026年时间维，无法构造$v_j(u_j)$；一张气候态速度图或空间四分位带也无法给出$a,b,\alpha$或释放时间区间。用10:30口岸冲击反推$a,b$属于事后标定，本文不采用。

在取得合格序列前，释放状态记为“未知”，不记为“未释放”；真实$t_r$、$\Delta_{r\to f}$、$t_a$及其预警尚不能估计。所需资料包括：发布前源区位移、裂隙开度、地震或声学释放信号，可核查的释放时刻和泥石流形成时刻，以及测量误差与缺测记录。取得这些资料后，才能把真实触发区间、真实预警和独立覆盖率作为主产品重新计算；本文不将当前结果表述为真实触发预报。

\subsection{一致性检查与适用限制}
本节的一致性检查包括：成员总数为3000，源区与口岸状态合计均为3000，支持宽度均为1 h，$p_{\mathrm{actionable}}$与状态质量一致，16个D8配置均提供有限首达，24种站点小时支持组合给出相同辅助区间。这些检查说明成员生成、支持传播和状态记录一致；融水代理不等同于真实释放或泥石流形成，区间也不具备真实频率覆盖。

主要不确定性如下。第一，IFS格点温度、短波和湿度是模式输出，跨高程与地表类型迁移可能产生系统偏差。第二，残差自回归、回归系数扰动和D8配置抽样组成模型集合不确定性，各分量未由多事件联合校准。第三，D8配置等权不是物理概率，改变配置集合会改变条件分位。第四，附件4速度场是气候态，只提供长期背景，缺少灾前加速序列；释放机制仍不可识别。第五，题面仅有一个口岸冲击时刻，无法检验长期覆盖率或误报率。因此，本节给出固定时刻的条件辅助产品及其适用范围，不将单事件结果升级为真实灾害确定性预测。
